% !TEX program = xelatex
% ================================================================
%  全国大学生电子设计竞赛 —— 设计报告 LaTeX 模板
%  参照 2024/2026 省赛报告格式要求 + 2025 电赛 G 题样例报告结构
%  编译方式：XeLaTeX → XeLaTeX（如需参考文献再跑一次 BibTeX）
%
%  格式要求摘要：
%    · 正文 A4 纸 8 页以内，纵向打印
%    · 首页另附 300 字以内中文摘要，无需英文翻译
%    · 正文小四号宋体（12pt），行距固定值 22 磅
%    · 标题字号自定（本模板：一级小三黑体、二级四号黑体、三级小四黑体）
%    · 每页上方留出 3 cm 以上空白，空白区域内不得有任何文字
%    · 每页右下端注明页码
% ================================================================

\documentclass[a4paper, zihao=-4, linespread=1.5, AutoFakeBold]{ctexart}
%  zihao=-4 → 小四号字（12pt）
%  linespread=1.5 配合下方 \setlength{\baselineskip}{22pt} 达到固定值 22 磅

% ====================== 宏包 ======================
\usepackage{amsmath, amssymb, amsfonts}
\usepackage{graphicx}
\usepackage{booktabs}
\usepackage{multirow}
\usepackage{array}
\usepackage{geometry}
\usepackage{setspace}
\usepackage{enumitem}
\usepackage{caption}
\usepackage{float}
\usepackage{fancyhdr}
\usepackage{titlesec}
\usepackage{listings}
\usepackage{xcolor}
\usepackage{hyperref}

% ====================== 页面布局 ======================
\geometry{
    a4paper,
    top      = 3.0cm,   % ≥3 cm 留白
    bottom   = 2.5cm,
    left     = 2.5cm,
    right    = 2.5cm,
    headsep  = 0.4cm,
    footskip = 1.0cm,
}

% ====================== 行距（固定值 22 磅） ======================
\setlength{\baselineskip}{22pt}
\setstretch{1.0}

% ====================== 中文字体 ======================
% ctexart 已自动配置 SimSun/SimHei/KaiTi/FangSong
% 如需手动指定，取消下方注释：
% \setCJKmainfont{SimSun}[BoldFont=SimHei, ItalicFont=KaiTi]
% \setCJKsansfont{SimHei}
% \setCJKmonofont{FangSong}
\setmainfont{Times New Roman}
\setsansfont{Arial}

% 华文中宋（封面大标题用）
\setCJKfamilyfont{zhongsong}{STZhongsong}[AutoFakeBold]
\newcommand{\zhongsong}{\CJKfamily{zhongsong}}

% ====================== 页眉页脚 ======================
\pagestyle{fancy}
\fancyhf{}
\fancyhead{}                            % 页眉留空
\fancyfoot[R]{\thepage}                 % 页码右下
\renewcommand{\headrulewidth}{0pt}
\renewcommand{\footrulewidth}{0pt}

% ====================== 章节标题格式 ======================
% 一级标题：小三号黑体（~15pt），中文编号"一、二、三……"
\ctexset{
    section = {
        format      = \fontsize{14pt}{18pt}\selectfont\bfseries\heiti,
        name        = {,、},           % "一、" "二、" ...
        number      = \chinese{section},
        aftername   = {},
        beforeskip  = 1.0ex plus 0.2ex minus 0.1ex,
        afterskip   = 0.5ex plus 0.1ex,
    },
    subsection = {
        format      = \large\bfseries\heiti,
        aftername   = {\quad},
        beforeskip  = 0.8ex plus 0.1ex,
        afterskip   = 0.3ex,
    },
    subsubsection = {
        format      = \normalsize\bfseries\heiti,
        aftername   = {\quad},
        beforeskip  = 0.5ex,
        afterskip   = 0.2ex,
    },
}

% ====================== 图表标题 ======================
\captionsetup{
    font          = {small},
    labelsep      = {space},
    justification = {centering},
}
\captionsetup[figure]{name={图}}
\captionsetup[table]{name={表}}

% ====================== 公式编号 ======================
\renewcommand{\theequation}{\arabic{section}-\arabic{equation}}

% ====================== 超链接 ======================
\hypersetup{
    colorlinks        = true,
    linkcolor         = black,
    citecolor         = black,
    urlcolor          = blue,
    bookmarksnumbered = true,
}

% ====================== 代码样式 ======================
\lstset{
    language        = C,
    basicstyle      = \small\ttfamily,
    keywordstyle    = \bfseries\color{blue!70!black},
    commentstyle    = \itshape\color{green!50!black},
    stringstyle     = \color{red!60!black},
    numbers         = left,
    numberstyle     = \tiny\color{gray},
    frame           = single,
    breaklines      = true,
    tabsize         = 4,
    showstringspaces = false,
}

% ====================== 段落 ======================
\setlength{\parindent}{2em}
\setlength{\parskip}{0pt}

% ====================== 自定义命令 ======================
\newcommand{\题号}[1]{\def\theProblemNum{#1}}
\newcommand{\参赛编号}[1]{\def\theTeamNum{#1}}
\newcommand{\队员}[1]{\def\theMembers{#1}}

% 封面页
\newcommand{\makereporttitle}{%
    \thispagestyle{empty}
    \begin{center}
        \vspace*{3cm}
        {\fontsize{22pt}{28pt}\selectfont\zhongsong\bfseries \theProblemNum}\\[2.0cm]
    \end{center}
}

% 图片占位框（开发阶段用，正式稿请替换为 \includegraphics）
\newcommand{\imgplaceholder}[2][0.7\textwidth]{%
    \fbox{\parbox{#1}{\centering\vspace{1.5cm}%
    {\kaishu #2}\\{\small （请替换为实际图片）}%
    \vspace{1.5cm}}}%
}

% ==================================================================
%                            正  文
% ==================================================================
\begin{document}

% ==================== 题目信息（← 请修改此处） ====================
\题号{X题}
\参赛编号{XXXX}
\队员{姓名1、姓名2、姓名3}

% ==================== 封面 ====================
\makereporttitle

% ==================== 摘要页（首页，≤300字） ====================
\thispagestyle{fancy}
\setcounter{page}{0}   % 摘要页不计入正文页码
\noindent{\fontsize{12pt}{18pt}\selectfont\songti\bfseries 摘\quad 要：}本设计实现了\textbf{XXXX系统}，能够\textbf{完成XXX功能}。
系统以 STM32XXXX 为核心控制器，集成 XXX 模块、XXX 模块以及 XXX 模块等。
通过\textbf{XXX方法}，实现了\textbf{XXX效果}。
经测试，本系统能够\textbf{满足题目各项要求}，各项性能指标均达到设计预期。

\vspace{0.5cm}
\noindent{\songti\bfseries 关键词：}关键词1；关键词2；关键词3；关键词4；关键词5
% ↑ 3–5 个关键词，中文分号分隔，无需英文翻译

\newpage

% ==================== 正文（从第 1 页起） ====================
\setcounter{page}{1}

% ----------------------------------------------------------------
%                     一、系统方案论证
% ----------------------------------------------------------------
\section{系统方案论证}

\subsection{方案比较与选择}

\subsubsection{XXX方案选择}

\textbf{方案一：}XXXX方案。描述该方案的实现方式及其优缺点。

\textbf{方案二：}XXXX方案。描述该方案的实现方式及其优缺点。

\textbf{方案选择：}方案一的优势在于 XXX，但存在 XXX 不足；方案二 XXX。
综合考虑 XXX 因素，故选择方案 X。

% ---- 更多方案比较 ----
\subsubsection{XXX方案选择}

\textbf{方案一：}XXXX。

\textbf{方案二：}XXXX。

\textbf{方案选择：}XXXX。

\subsubsection{XXX方案选择}

\textbf{方案一：}XXXX。

\textbf{方案二：}XXXX。

\textbf{方案选择：}XXXX。

\subsection{总体方案描述}

系统总体框图如图~\ref{fig:system_block}~所示。
该系统以 STM32XXXX 为核心控制单元……
（描述总体架构，各模块功能及数据流向）

\begin{figure}[H]
    \centering
    % \includegraphics[width=0.8\textwidth]{figures/system_block.png}
    \imgplaceholder{系统总体框图}
    \caption{系统总体框图}
    \label{fig:system_block}
\end{figure}

% ----------------------------------------------------------------
%                    二、理论分析与计算
% ----------------------------------------------------------------
\section{理论分析与计算}

\subsection{XXX电路分析}

XXX电路的电压传递函数为：
\begin{equation}
    H(s) = \frac{V_{\text{out}}(s)}{V_{\text{in}}(s)}
         = \frac{a_0}{b_2 s^2 + b_1 s + b_0}
    \label{eq:transfer_func}
\end{equation}

经过分析与计算，各元件参数确定如下：
\begin{equation}
    R_1 = XXX\;\Omega, \quad
    C_1 = XXX\;\text{F}, \quad \ldots
    \label{eq:params}
\end{equation}

将式~(\ref{eq:transfer_func})与式~(\ref{eq:params})的分母不同阶次项系数对应，
可以确定电路中各电容与电阻的具体参数。

\begin{figure}[H]
    \centering
    % \includegraphics[width=0.7\textwidth]{figures/bode_plot.png}
    \imgplaceholder{波特图 / 频率响应图}
    \caption{XXX 频率响应特性（波特图）}
    \label{fig:bode}
\end{figure}

\subsection{XXX分析方法}

根据 XXX 原理，可以得到……

\begin{equation}
    f(x) = \sum_{n=0}^{\infty} a_n x^n
\end{equation}

% ----------------------------------------------------------------
%                   三、电路与程序设计
% ----------------------------------------------------------------
\section{电路与程序设计}

\subsection{XXX电路设计}

XXX电路如图~\ref{fig:circuit1}~所示，使用 XXX 芯片实现 XXX 功能。
主要元件及其作用如下：
\begin{itemize}[leftmargin=2em, itemsep=2pt]
    \item \textbf{XXX 芯片}：用于 XXX
    \item $R_1$、$C_1$：构成 XXX 电路
    \item 运算放大器 OPAxxxx：用于 XXX
\end{itemize}

\begin{figure}[H]
    \centering
    % \includegraphics[width=0.8\textwidth]{figures/circuit1.png}
    \imgplaceholder{电路原理图}
    \caption{XXX 电路原理图}
    \label{fig:circuit1}
\end{figure}

\subsection{XXX电路设计}

描述另一个电路模块的功能与设计……

\subsection{软件程序设计}

本设计采用 STM32XXXX 作为主控芯片，完成数据的采样与处理、
信号产生与控制等工作。程序设计流程图如图~\ref{fig:flowchart}~所示。

\begin{figure}[H]
    \centering
    % \includegraphics[width=0.45\textwidth]{figures/flowchart.png}
    \imgplaceholder[0.45\textwidth]{程序流程图}
    \caption{程序设计流程图}
    \label{fig:flowchart}
\end{figure}

系统软件主要包含以下功能模块：
\begin{enumerate}[leftmargin=2em, itemsep=2pt]
    \item \textbf{初始化模块}：完成系统时钟、GPIO、ADC、DAC、SPI 等外设配置；
    \item \textbf{信号采集模块}：通过 ADC + DMA 实现连续采样；
    \item \textbf{数据处理模块}：实现 XXX 算法（如 IIR 滤波、FFT 等）；
    \item \textbf{信号输出模块}：通过 DAC / DDS 输出目标信号；
    \item \textbf{人机交互模块}：按键输入与 OLED / LCD 显示。
\end{enumerate}

% ----------------------------------------------------------------
%                  四、测试方案与测试结果
% ----------------------------------------------------------------
\section{测试方案与测试结果}

\subsection{测试条件}

主要测试仪器设备如表~\ref{tab:instruments}~所示。

\begin{table}[H]
    \centering
    \caption{测试仪器设备清单}
    \label{tab:instruments}
    \begin{tabular}{ccc}
        \toprule
        \textbf{仪器名称} & \textbf{型号} & \textbf{数量} \\
        \midrule
        数字万用表   & XXXX & 一台 \\
        示波器       & XXXX & 一台 \\
        信号发生器   & XXXX & 一台 \\
        可调电源     & XXXX & 一台 \\
        \bottomrule
    \end{tabular}
\end{table}

\subsection{测试方案与测试结果}

\subsubsection{XXX测试}

测试方法：XXXX。测试结果如表~\ref{tab:test1}~所示。

\begin{table}[H]
    \centering
    \caption{XXX 测试记录表}
    \label{tab:test1}
    \begin{tabular}{cccc}
        \toprule
        \textbf{测试条件} & \textbf{理论值} & \textbf{实测值} & \textbf{相对误差} \\
        \midrule
        XXX & XXX & XXX & XXX \\
        XXX & XXX & XXX & XXX \\
        XXX & XXX & XXX & XXX \\
        XXX & XXX & XXX & XXX \\
        \bottomrule
    \end{tabular}
\end{table}

\subsubsection{XXX测试}

测试方法：XXXX。测试结果如表~\ref{tab:test2}~所示。

\begin{table}[H]
    \centering
    \caption{XXX 测试记录表}
    \label{tab:test2}
    \begin{tabular}{cccc}
        \toprule
        \textbf{测试条件} & \textbf{设定值} & \textbf{实际值} & \textbf{相对误差} \\
        \midrule
        XXX & XXX & XXX & XXX \\
        XXX & XXX & XXX & XXX \\
        XXX & XXX & XXX & XXX \\
        \bottomrule
    \end{tabular}
\end{table}

% 可选：插入测试波形截图
% \begin{figure}[H]
%     \centering
%     \includegraphics[width=0.7\textwidth]{figures/test_waveform.png}
%     \caption{XXX 测试波形}
%     \label{fig:test_wave}
% \end{figure}

\subsection{测试结果分析}

根据以上测试结果，各项实验参数误差均在 X\% 之内，
\textbf{达到了设计要求}。
本实验的主要误差来源于 XXX 以及 XXX。

为进一步提高系统性能，可以从以下方面改进：
\begin{enumerate}[leftmargin=2em, itemsep=2pt]
    \item XXX；
    \item XXX；
    \item XXX。
\end{enumerate}

% ----------------------------------------------------------------
%                       五、参考文献
% ----------------------------------------------------------------
\section{参考文献}

% 不要写"无"！至少列 3–5 条与题目相关的参考文献。
\begin{thebibliography}{99}

\bibitem{ref1}
    作者姓名.
    \textit{书名}[M].
    出版地：出版社，年份.

\bibitem{ref2}
    作者1，作者2.
    \textit{书名——副标题}[M].
    出版地：出版社，年份.

\bibitem{ref3}
    作者.
    论文题目[J].
    \textit{期刊名}，年份，卷(期)：起始页--终止页.

\bibitem{ref4}
    公司名.
    \textit{芯片型号 Datasheet}[Z].
    年份.

\end{thebibliography}

% ==================== 附录（可选，不计入页数） ====================
\newpage
\appendix
\renewcommand{\thesection}{附录\Alph{section}}

\section{核心代码片段}

\begin{lstlisting}[
    caption = {XXX 核心代码},
    label   = {lst:code1},
]
#include "stm32g4xx_hal.h"

// 在此粘贴关键代码
void SystemClock_Config(void) {
    // 系统时钟配置
}

void ADC_Init(void) {
    // ADC 初始化
}
\end{lstlisting}

\section{完整电路原理图}

\begin{figure}[H]
    \centering
    % \includegraphics[width=\textwidth]{figures/full_schematic.png}
    \imgplaceholder[\textwidth]{完整电路原理图}
    \caption{完整电路原理图}
\end{figure}

\section{实物照片}

\begin{figure}[H]
    \centering
    % \includegraphics[width=0.7\textwidth]{figures/photo.jpg}
    \imgplaceholder{系统实物照片}
    \caption{系统实物照片}
\end{figure}

\end{document}
